%% GENERATED by python/paper/build.py from results/ -- do not edit by hand.
%% Rebuild:  .venv\Scripts\python.exe python\paper\build.py --only OECD26_Omega_nuExit
%% Source:   results/oecd26/layer2.csv, anchor.csv; data/oecdFigure1.csv, oecd26_theta.csv
\begin{table}[!htb]
\centering
\begin{threeparttable}
\caption{The countries of the cross-country test, each calibrated to its pension spending, $\nu$ corrected at the effective exit age}
\label{table:OECD26:omega_nuExit}
\renewcommand{\arraystretch}{1.25}
\footnotesize\setlength{\tabcolsep}{3pt}
\begin{tabularx}{\textwidth}{lYYYYYYYY}
\toprule
 & \multicolumn{5}{c}{\textbf{Data}} & \multicolumn{3}{c}{\textbf{Model, own $\omega$}} \\
\cmidrule(lr){2-6}\cmidrule(lr){7-9}
\textbf{Country} & $\nu_{2020}$ & $\tilde V$ & spending & $I$ & $\theta$ & $\omega$ & $\theta$ & $I$ \\
\midrule 
Austria & 0.95 & 0.161 & 14.3 & 0.791 & 0.736 & 1.315 & 0.830 & 0.854 \\ % row: aut
Belgium & 0.87 & 0.104 & 11.5 & 0.616 & 0.376 & 1.106 & 0.805 & 0.837 \\ % row: bel
Canada & 1.28 & 0.171 & 6.4 & 0.595 & 0.318 & 1.144 & 0.228 & 0.564 \\ % row: can
Czech Republic & 0.91 & 0.170 & 8.6 & 0.603 & 0.342 & 0.966 & 0.632 & 0.731 \\ % row: cze
Finland & 0.97 & 0.094 & 12.6 & 0.944 & 0.941 & 1.228 & 0.793 & 0.829 \\ % row: fin
France & 0.84 & 0.150 & 14.4 & 0.849 & 0.822 & 1.156 & 0.828 & 0.853 \\ % row: fra
Germany & 0.93 & 0.148 & 11.2 & 0.778 & 0.714 & 1.097 & 0.732 & 0.789 \\ % row: deu
Greece & 0.79 & 0.225 & 17.9 & 0.857 & 0.833 & 1.305 & 0.874 & 0.888 \\ % row: grc
Hungary & 0.78 & 0.265 & 7.7 & 0.997 & 0.997 & 0.788 & 0.572 & 0.700 \\ % row: hun
Ireland & 1.28 & 0.240 & 3.3 & 0.500 & 0.000 & 0.876 & 0.000 & 0.500 \\ % row: irl
Italy & 0.88 & 0.134 & 17.4 & 0.956 & 0.954 & 1.406 & 0.893 & 0.903 \\ % row: ita
Japan & 1.27 & 0.124 & 9.6 & 0.747 & 0.662 & 1.387 & 0.650 & 0.740 \\ % row: jpn
Luxembourg & 1.02 & 0.187 & 9.0 & 0.848 & 0.820 & 1.068 & 0.608 & 0.718 \\ % row: lux
New Zealand & 1.80 & 0.141 & 5.1 & 0.501 & 0.003 & 1.504 & 0.000 & 0.500 \\ % row: nzl
Norway & 1.27 & 0.101 & 7.6 & 0.593 & 0.313 & 1.253 & 0.476 & 0.656 \\ % row: nor
Poland & 0.84 & 0.240 & 11.4 & 0.964 & 0.962 & 1.002 & 0.751 & 0.801 \\ % row: pol
Portugal & 1.04 & 0.249 & 13.6 & 0.892 & 0.879 & 1.370 & 0.808 & 0.839 \\ % row: prt
Spain & 0.92 & 0.183 & 12.8 & 0.877 & 0.859 & 1.168 & 0.791 & 0.827 \\ % row: esp
Sweden & 1.28 & 0.081 & 8.5 & 0.619 & 0.383 & 1.329 & 0.571 & 0.700 \\ % row: swe
United Kingdom & 1.07 & 0.185 & 7.8 & 0.550 & 0.181 & 1.116 & 0.623 & 0.726 \\ % row: gbr
United States$^{\ast}$ & 1.34 & 0.333 & 7.4 & 0.687 & 0.544 & 1.272 & 0.544 & 0.687 \\ % row: usa
\bottomrule
\end{tabularx}
\begin{tablenotes}[flushleft]
\footnotesize
\item[] \textit{Note:} 21 countries. $\nu_{2020}$, workers per retiree; $\tilde V$, the dispersion of relative incomes at the U.S.\ income percentiles; spending in \% of GDP; $\theta$, the earnings-linked share of the least-squares line through the public pension at 0.5, 1 and 2 times average earnings ($c/(a+c)$ of $p(y) = a + c\,y$), times one minus the share of public spending on minimum top-ups and tested benefits, and $I = 1/(2-\theta)$, the Bismarckian index. $\omega$ is calibrated to spending at the observed design at $\beta = 0.65$ and the U.S.\ $\kappa = 93.71$, $\lambda_{n,t} = 2\kappa\,\pi_{n,t}(1-\pi_{n,t})$, $\pi_{n,t} = \Pi_{n,t}/R_{n,t}$; the last two columns are the model\textquotesingle s design, in force in 2020 on the freely simulated path, chosen in 1990, and its index. $^{\ast}$The U.S.\ is the anchor, where $\omega$ and $\kappa$ are calibrated. Turnout: each country\textquotesingle s own propensities to vote by income group, from the CSES; Luxembourg, which has no CSES survey and enforces compulsory voting, at a flat profile.
\end{tablenotes}
\end{threeparttable}
\end{table}
